\begin{lemma}
\label{lemma-split-off-proper-part-henselian}
\begin{reference}
A reference for the case of an adic Noetherian base is
\cite[III, Proposition 5.5.1]{EGA}
\end{reference}
Let $(A, I)$ be a henselian pair. Let $X \to \Spec(A)$
be separated and of finite type. Set $X_0 = X \times_{\Spec(A)} \Spec(A/I)$.
Let $Y \subset X_0$ be an open and closed subscheme such that
$Y \to \Spec(A/I)$ is proper. Then there exists an open and closed
subscheme $W \subset X$ which is proper over $A$ with
$W \times_{\Spec(A)} \Spec(A/I) = Y$.
\end{lemma}

\begin{proof}
We will denote $T \mapsto T_0$ the base change by $\Spec(A/I) \to \Spec(A)$.
By Chow's lemma (in the form of
Limits, Lemma \ref{limits-lemma-chow-finite-type})
there exists a surjective proper morphism $\varphi : X' \to X$ such
that $X'$ admits an immersion into $\mathbf{P}^n_A$.
Set $Y' = \varphi^{-1}(Y)$. This is an open and closed subscheme
of $X'_0$. Suppose the lemma holds for $(X', Y')$. Let $W' \subset X'$
be the open and closed subscheme proper over $A$ such that $Y' = W'_0$.
By Morphisms, Lemma \ref{morphisms-lemma-image-proper-scheme-closed}
$W = \varphi(W') \subset X$ and
$Q = \varphi(X' \setminus W') \subset X$ are closed subsets and by
Morphisms, Lemma \ref{morphisms-lemma-image-proper-is-proper}
$W$ is proper over $A$. The image of $W \cap Q$ in $\Spec(A)$ is closed.
Since $(A, I)$ is henselian, if $W \cap Q$ is nonempty, then we
find that $W \cap Q$ has a point lying over $\Spec(A/I)$.
This is impossible as $W'_0 = Y' = \varphi^{-1}(Y)$.
We conclude that $W$ is an open and closed subscheme
of $X$ proper over $A$ with $W_0 = Y$.
Thus we reduce to the case described in the next paragraph.

\medskip\noindent
Assume there exists an immersion $j : X \to \mathbf{P}^n_A$ over $A$.
Let $\overline{X}$ be the scheme theoretic image of $j$.
Since $j$ is a quasi-compact morphism
(Schemes, Lemma \ref{schemes-lemma-quasi-compact-permanence})
we see that $j : X \to \overline{X}$ is an open immersion
(Morphisms, Lemma \ref{morphisms-lemma-quasi-compact-immersion}).
Hence the base change $j_0 : X_0 \to \overline{X}_0$
is an open immersion as well.
Thus $j_0(Y) \subset \overline{X}_0$ is open.
It is also closed by Morphisms, Lemma
\ref{morphisms-lemma-image-proper-scheme-closed}.
Suppose that the lemma holds for $(\overline{X}, j_0(Y))$.
Let $\overline{W} \subset \overline{X}$ be the
corresponding open and closed subscheme proper over $A$
such that $j_0(Y) = \overline{W}_0$.
Then $T = \overline{W} \setminus j(X)$ is closed in $\overline{W}$,
hence has closed image in $\Spec(A)$ by properness of $\overline{W}$
over $A$. Since $(A, I)$ is henselian, we find that if $T$
is nonempty, then there is a point of $T$ mapping into $\Spec(A/I)$.
This is impossible because $j_0(Y) = \overline{W}_0$ is contained in $j(X)$.
Hence $\overline{W}$ is contained in $j(X)$ and we can
set $W \subset X$ equal to the unique open and closed
subscheme mapping isomorphically to $\overline{W}$ via $j$.
Thus we reduce to the case described in the next paragraph.

\medskip\noindent
Assume $X \subset \mathbf{P}^n_A$ is a closed subscheme.
Then $X \to \Spec(A)$ is a proper morphism.
Let $Z = X_0 \setminus Y$. This is an open and closed
subscheme of $X_0$ and $X_0 = Y \amalg Z$.
Let $X \to X' \to \Spec(A)$ be the Stein factorization as in
Theorem \ref{theorem-stein-factorization-general}.
Let $Y' \subset X'_0$ and $Z' \subset X'_0$ be the images of
$Y$ and $Z$.
Since the fibres of $X \to Z$ are geometrically connected,
we see that $Y' \cap Z' = \emptyset$.
Hence $X'_0 = Y' \amalg Z'$ as $X \to X'$ is surjective.
Since $X' \to \Spec(A)$ is integral, we see that
$X'$ is the spectrum of an $A$-algebra integral over $A$.
Recall that open and closed subsets of spectra correspond
$1$-to-$1$ with idempotents in the corresponding ring, see
Algebra, Lemma \ref{algebra-lemma-disjoint-decomposition}.
Hence by
More on Algebra, Lemma \ref{more-algebra-lemma-characterize-henselian-pair}
we see that we may write $X' = W' \amalg V'$
with $W'$ and $V'$ open and closed and
with $Y' = W'_0$ and $Z' = V'_0$.
Let $W$ be the inverse image in $X$
to finish the proof.
\end{proof}